\documentclass[10pt]{article}
\usepackage[letterpaper,landscape,margin=0.55in]{geometry}
\usepackage{graphicx}
\usepackage{booktabs}
\usepackage{amsmath}
\usepackage{float}
\usepackage[hidelinks]{hyperref}
\setlength{\parindent}{0pt}
\title{Q-Weighted Politician Event Studies}
\author{Exploratory fixed-effect comparison}
\date{August 18, 2026}
\begin{document}
\maketitle
\textbf{Connection to the paper.} This exploratory analysis applies the 
corrective Q-weighted stacked difference-in-differences estimator from 
\href{https://doi.org/10.3386/w32054}{Wing, Freedman, and Hollingsworth 
(2024), \textit{Stacked Difference-in-Differences}}. We replicate their 
method and their original fixed-effect specification, not their numerical 
Medicaid results: the estimates below are for fires in this project. The 
paper's preferred fixed-composition target is a trimmed aggregate ATT.

\textbf{What cohort ID means.} A cohort ID labels one stacked 
sub-experiment (a province-election pair); it does not label a control type. 
Every included cohort contains both treated observations ($treat=1$) and 
controls ($treat=0$). All $treat=0$ observations are included because they 
remain untreated throughout that stack's analysis window; no distinction 
among control types is used.

\textbf{All cohorts are included.} No cohort is discarded for lacking 
some leads or lags. Each sub-experiment contributes only at the event years 
available in the data. Thus IDs 3--6 contribute from -5 through 4; IDs 1--2 
enter later; and IDs 7--8 leave after event year 2.
\begin{center}
\small
\begin{tabular}{rllll}
\toprule
ID & Province & Election & Available event years & Used? \\
\midrule
1 & Haryana & 2014-11 & -3 to 4 & Include \\
2 & Bihar & 2015-12 & -4 to 4 & Include \\
3 & Punjab & 2017-04 & -5 to 4 & Include \\
4 & Uttar Pradesh & 2017-04 & -5 to 4 & Include \\
5 & Haryana & 2019-11 & -5 to 4 & Include \\
6 & Bihar & 2020-12 & -5 to 4 & Include \\
7 & Punjab & 2022-04 & -5 to 2 & Include \\
8 & Uttar Pradesh & 2022-04 & -5 to 2 & Include \\
\bottomrule
\end{tabular}
\end{center}
\textbf{Corrective weights.} For cohort $s$ and event year $e$, treated 
observations receive weight one and controls receive
\[Q_{se}=\frac{n^T_{se}/N^T_e}{n^C_{se}/N^C_e}.\]
At each event year, this makes the weighted cohort shares among controls 
match the shares among treated observations for the cohorts available at 
that event year, which is the paper's corrective-weight principle.

\textbf{Estimand caveat.} Because cohort availability changes over event 
time, this requested all-cohort specification is not the paper's 
fixed-composition trimmed aggregate ATT over -5 through 4. Event year -5 
uses IDs 3--8; -4 uses IDs 2--8; -3 through 2 use IDs 1--8; and 3--4 use 
IDs 1--6. Consequently, differences across event years can reflect both 
dynamic effects and changes in the contributing cohorts.
\clearpage
\section*{Event-study estimates}
Event years run from -5 through 4, with -1 omitted. 
The sample contains 12,599,763 observations from cohort IDs 1, 2, 3, 4, 5, 6, 7, 8. All models include wind direction and average wind speed, use Q-weights, and cluster standard errors by AC $\times$ cohort.
 Control weights range from 0.167 to 1.938.

\textbf{Specifications.} Original stacked FE absorbs treatment-group and 
relative-year effects. FE 01 absorbs grid $\times$ cohort effects. FE 05 
absorbs grid $\times$ cohort and province $\times$ cohort $\times$ 
election-year effects.

\begin{figure}[H]
\centering
\begin{minipage}[t]{0.325\textwidth}
\centering
\textbf{Original stacked FE}\par\smallskip
\includegraphics[width=\linewidth]{\detokenize{C:/Users/eunic/OneDrive/Documents/GitHub/proj_bureaucrats_farms/figures/exploratory_analysis/politician_byprov_stack_weights/politician_qweight_original_ori.png}}
\end{minipage}
\hfill
\begin{minipage}[t]{0.325\textwidth}
\centering
\textbf{FE 01}\par\smallskip
\includegraphics[width=\linewidth]{\detokenize{C:/Users/eunic/OneDrive/Documents/GitHub/proj_bureaucrats_farms/figures/exploratory_analysis/politician_byprov_stack_weights/politician_qweight_fe01_ori.png}}
\end{minipage}
\hfill
\begin{minipage}[t]{0.325\textwidth}
\centering
\textbf{FE 05}\par\smallskip
\includegraphics[width=\linewidth]{\detokenize{C:/Users/eunic/OneDrive/Documents/GitHub/proj_bureaucrats_farms/figures/exploratory_analysis/politician_byprov_stack_weights/politician_qweight_fe05_ori.png}}
\end{minipage}
\caption{Q-weighted event-study estimates and 95\% confidence intervals.}
\end{figure}
\clearpage
\section*{Q-weighted DiD interactions}
The two post-treatment estimates are intentionally offset horizontally. 
Their coefficient values and confidence intervals are unchanged; the added 
space only makes overlapping 90\% and 95\% intervals visible.
\begin{figure}[H]
\centering
\begin{minipage}[t]{0.44\textwidth}
\centering
\textbf{Original stacked FE}\par\smallskip
\includegraphics[width=\linewidth]{\detokenize{C:/Users/eunic/OneDrive/Documents/GitHub/proj_bureaucrats_farms/figures/exploratory_analysis/politician_byprov_stack_weights/politician_qweight_original_interaction.png}}
\end{minipage}
\hfill
\begin{minipage}[t]{0.44\textwidth}
\centering
\textbf{FE 01}\par\smallskip
\includegraphics[width=\linewidth]{\detokenize{C:/Users/eunic/OneDrive/Documents/GitHub/proj_bureaucrats_farms/figures/exploratory_analysis/politician_byprov_stack_weights/politician_qweight_fe01_interaction.png}}
\end{minipage}
\par\medskip
\begin{minipage}[t]{0.44\textwidth}
\centering
\textbf{FE 05}\par\smallskip
\includegraphics[width=\linewidth]{\detokenize{C:/Users/eunic/OneDrive/Documents/GitHub/proj_bureaucrats_farms/figures/exploratory_analysis/politician_byprov_stack_weights/politician_qweight_fe05_interaction.png}}
\end{minipage}
\caption{Q-weighted pre/post interactions with separated post positions.}
\end{figure}
\clearpage
\section*{Regression estimates}

\begingroup
\centering
\begin{tabular}{lccc}
   \tabularnewline \midrule \midrule
   Dependent Variable: & \multicolumn{3}{c}{countk}\\
   Specification                          & Original stacked FE & FE 01        & FE 05 \\   
   Model:                                 & (1)                 & (2)          & (3)\\  
   \midrule
   \emph{Variables}\\
   treat $\times$ relative\_year $=$ -5   & -4.52               & 3.58         & 3.04\\   
                                          & (9.67)              & (5.24)       & (5.28)\\   
   treat $\times$ relative\_year $=$ -4   & -3.80               & -10.6$^{**}$ & -11.2$^{**}$\\   
                                          & (7.22)              & (4.63)       & (4.76)\\   
   treat $\times$ relative\_year $=$ -3   & -3.20               & 4.99         & 4.93\\   
                                          & (6.04)              & (5.38)       & (5.54)\\   
   treat $\times$ relative\_year $=$ -2   & 1.48                & -5.55        & -5.55\\   
                                          & (5.76)              & (4.96)       & (4.96)\\   
   treat $\times$ relative\_year $=$ 0    & 5.51                & 2.37         & 3.94\\   
                                          & (7.05)              & (4.73)       & (4.56)\\   
   treat $\times$ relative\_year $=$ 1    & 7.25                & -4.67        & -3.10\\   
                                          & (7.02)              & (4.66)       & (4.58)\\   
   treat $\times$ relative\_year $=$ 2    & 8.35                & 12.3$^{**}$  & 13.8$^{**}$\\   
                                          & (8.62)              & (6.16)       & (5.89)\\   
   treat $\times$ relative\_year $=$ 3    & 7.02                & 6.13         & 7.53\\   
                                          & (9.78)              & (5.34)       & (5.96)\\   
   treat $\times$ relative\_year $=$ 4    & 3.85                & 8.35         & 10.8$^{*}$\\   
                                          & (10.0)              & (5.84)       & (6.25)\\   
   \midrule
   \emph{Fixed-effects}\\
   treat                                  & Yes                 &              & \\  
   relative\_year                         & Yes                 &              & \\  
   grid\_cohort                           &                     & Yes          & Yes\\  
   province-cohort\_id-election\_year     &                     &              & Yes\\  
   \midrule
   \emph{Fit statistics}\\
   Observations                           & 12,599,763          & 12,599,763   & 12,599,763\\  
   R$^2$                                  & 0.00456             & 0.12969      & 0.13015\\  
   Within R$^2$                           & 0.00441             & 0.00487      & 0.00485\\  
   \midrule \midrule
   \multicolumn{4}{l}{\emph{Clustered (cluster\_id) standard-errors in parentheses}}\\
   \multicolumn{4}{l}{\emph{Signif. Codes: ***: 0.01, **: 0.05, *: 0.1}}\\
\end{tabular}
\par\endgroup


}
\vfill
\footnotesize\textit{Notes:} The omitted event year is -1. The dependent 
variable is the project-standard fire count multiplied by 1,000. The same 
complete-case, all-cohort sample and Q-weights are used in all three 
columns. Cohort composition varies with event-time coverage. Method: 
\href{https://doi.org/10.3386/w32054}{Wing, 
Freedman, and Hollingsworth (2024)}. Implementation reference: 
\href{https://github.com/hollina/stacked-did-weights}{authors' example 
code}.
\end{document}
